% answers_short.tex —— 2014 年·四川卷·文科 · 简短答案速查
% 由 tools/make_short_answers.py 自 answers.tex 的 \gksolhead/\gkanswer 提取，
% 并与 2014/2014_数学真题_四川文_答案要点.md 逐题比对后生成，勿手改（改源 MD 后重跑）。

\begin{gknotice}
\textbf{使用说明}：本节只给最终答案便于快速核对，完整推导见解析版。答案为编写者独立推导，\textbf{非官方标准答案}。
\end{gknotice}

\gkpart{选择题}
\begin{center}\begin{tabular}{ccccccccccc}
\toprule
题号 & 1 & 2 & 3 & 4 & 5 & 6 & 7 & 8 & 9 & 10 \\
\midrule
答案 & D & A & A & D & D & C & B & C & B & B \\
\bottomrule
\end{tabular}\end{center}

\gkpart{填空题}
\begin{itemize}[leftmargin=2.6em]
  \item[\textbf{11.}] $\dfrac{\sqrt5}2$。$a=2$，$b=1$，$c=\sqrt{a^2+b^2}=\sqrt5$，故 $e=\dfrac ca=\dfrac{\sqrt5}2$。
  \item[\textbf{12.}] $-2\mathrm i$。分子分母同乘 $1-\mathrm i$，得 $\dfrac{(2-2\mathrm i)(1-\mathrm i)}2=\dfrac{-4\mathrm i}2=-2\mathrm i$。
  \item[\textbf{13.}] $1$。周期为 $2$，故 $f\left(\dfrac32\right)=f\left(-\dfrac12\right)=-4\times\dfrac14+2=1$。
  \item[\textbf{14.}] $2$。$\boldsymbol c=(m+4,2m+2)$，由两夹角相等得 $\dfrac{5m+8}{\sqrt5}=\dfrac{8m+20}{2\sqrt5}$，解得 $m=2$。
  \item[\textbf{15.}] ①③④。①为满射的等价表述；②错，如 $\arctan x$ 有界但无最值；③对，否则 $f=(f+g)-g$ 也有界；④对，$a>0$ 或 $a<0$ 时 $f$ 无最大值，只能 $a=0$，此时 $f=\dfrac{x}{x^2+1}$ 有最大值和最小值。
\end{itemize}

\gkpart{解答题}
\begin{itemize}[leftmargin=2.6em]
  \item[\textbf{16.}] (1) $\dfrac19$。有放回抽取 $3$ 次共 $27$ 个等可能结果，满足 $a+b=c$ 的只有 $(1,1,2)$，$(1,2,3)$，$(2,1,3)$ 三个；(2) $\dfrac89$。$a$，$b$，$c$ 完全相同的只有 $(1,1,1)$，$(2,2,2)$，$(3,3,3)$ 三个，故所求为 $1-\dfrac3{27}=\dfrac89$
  \item[\textbf{17.}] (1) $\left[-\dfrac\pi4+\dfrac{2k\pi}3,\ \dfrac\pi{12}+\dfrac{2k\pi}3\right]$（$k\in\mathbf Z$）。由 $-\dfrac\pi2+2k\pi\leqslant 3x+\dfrac\pi4\leqslant\dfrac\pi2+2k\pi$ 解出；(2) $-\sqrt2$ 或 $-\dfrac{\sqrt5}2$。记 $u=\cos\alpha-\sin\alpha$、$v=\cos\alpha+\sin\alpha$，条件化为 $v=\dfrac45u^2v$：$v=0$ 给 $\alpha=\dfrac{3\pi}4$、$u=-\sqrt2$；$u^2=\dfrac54$ 给 $u=-\dfrac{\sqrt5}2$（$\alpha\approx97.2^\circ$ 或 $172.8^\circ$），两支均满足第二象限条件
  \item[\textbf{18.}] (1) 证明。矩形 $ACC_1A_1$ 给出 $CC_1\perp AC$，又 $BC\perp AC$，且 $AC\cap CC_1=C$，故 $BC\perp$ 平面 $ACC_1A_1$；(2) 存在，$M$ 为 $AB$ 的中点。以 $\overrightarrow{AB}$，$\overrightarrow{AC}$，$\overrightarrow{AA_1}$ 为基表示，$M$ 为中点时 $\overrightarrow{DE}=-\overrightarrow{A_1M}+\dfrac12\overrightarrow{A_1C}$，故 $DE\parallel$ 平面 $A_1MC$；由系数比较知这样的 $M$ 唯一
  \item[\textbf{19.}] (1) 证明。$b_n=2^{a_n}>0$，且 $\dfrac{b_{n+1}}{b_n}=2^{a_{n+1}-a_n}=2^d$ 为常数，故 $\{b_n\}$ 是公比为 $2^d$ 的等比数列；(2) $S_n=\dfrac{(3n-1)4^{n+1}+4}9$。切线截距 $a_2-\dfrac1{\ln2}=2-\dfrac1{\ln2}$ 给出 $a_2=2$，故 $a_n=n$，$b_n=2^n$，$a_nb_n^2=n\cdot4^n$，错位相减即得
  \item[\textbf{20.}] (1) $\dfrac{x^2}6+\dfrac{y^2}2=1$。$c=2$，由 $e=\dfrac{\sqrt6}3$ 得 $a=\sqrt6$，$b^2=a^2-c^2=2$；(2) $2\sqrt3$。设 $T(-3,m)$，过 $F$ 的垂线为 $x=-2+ms$，$y=s$，代入椭圆得 $(m^2+3)s^2-4ms-2=0$；由 $P+Q=T$ 得 $m^2=1$，面积 $=|x_Py_Q-x_Qy_P|=2|s_1-s_2|=2\sqrt3$
  \item[\textbf{21.}] (1) 当 $a\leqslant\dfrac12$ 时 $g_{\min}=1-b$；当 $\dfrac12<a<\dfrac{\mathrm e}2$ 时 $g_{\min}=2a-2a\ln(2a)-b$；当 $a\geqslant\dfrac{\mathrm e}2$ 时 $g_{\min}=\mathrm e-2a-b$。由 $g'(x)=\mathrm e^x-2a$ 的零点位置分三种情形求得；(2) 证明。由 $f(1)=0$ 得 $b=\mathrm e-1-a$，且 $f(0)=0$；对 $[0,x_0]$ 与 $[x_0,1]$ 用罗尔定理知 $g=f'$ 在 $(0,1)$ 内有两个零点，而 $g''=\mathrm e^x>0$ 使 $g$ 严格凸，故 $g(0)>0$、$g(1)>0$，即 $a+2-\mathrm e>0$、$1-a>0$，从而 $\mathrm e-2<a<1$
\end{itemize}
